\documentclass[../../../uem_phy_std_q.tex]{subfiles}

\begin{document}
    \setcounter{figure}{0}
    抵抗A（抵抗値$10\:\Omega$）と抵抗B（抵抗値$40\:\Omega$），%
    電池（起電力$10\:\mathrm{V}$，内部抵抗無視），および図\ref{fig:fig_uem_phy_std_em_cx_02_01_q}に示すような%
    電流・電圧特性をもつ豆電球を用意する．

    \begin{enumerate}
        \item 豆電球を電池に接続する．豆電球を流れる電流の大きさを求めよ．
        \item 同じ豆電球を2つ直列につなぎ，電池に接続する．豆電球を流れる電流の大きさを求めよ．
        \item 図\ref{fig:fig_uem_phy_std_em_cx_02_02_q}のように抵抗Aと豆電球を直列につなぎ，電池に接続する．%
            豆電球を流れる電流の大きさを求めよ．
        \item 図\ref{fig:fig_uem_phy_std_em_cx_02_03_q}のように抵抗Bを豆電球に並列につなぎ，%
            これらに抵抗Aを直列につないで，%
            電池に接続する．豆電球を流れる電流の大きさを求めよ．
    \end{enumerate}
    \begin{minipage}{\linewidth}
        \centering
            \includegraphics{\subfix{fig_uem_phy_std_em_cx_02/fig_uem_phy_std_em_cx_02_01_q.pdf}}
            \captionof{figure}{} \label{fig:fig_uem_phy_std_em_cx_02_01_q}
            \vspace{5pt}            
            \includegraphics{\subfix{fig_uem_phy_std_em_cx_02/fig_uem_phy_std_em_cx_02_02_q.pdf}}
            \captionof{figure}{} \label{fig:fig_uem_phy_std_em_cx_02_02_q}
            \vspace{5pt}
            \includegraphics{\subfix{fig_uem_phy_std_em_cx_02/fig_uem_phy_std_em_cx_02_03_q.pdf}}
            \captionof{figure}{} \label{fig:fig_uem_phy_std_em_cx_02_03_q}
    \end{minipage}        
\end{document}






